\documentclass[article,12pt]{elsarticle}
\usepackage[spanish]{babel}
\usepackage[utf8]{inputenc}
\usepackage{amsmath, amssymb, graphicx, float, caption, subfigure, longtable, lmodern, booktabs}

\begin{document}

\begin{frontmatter}
    \title{Bitácora: Laboratorio de Electromagnetismo \\ Práctica 1: Electrostática}
    \author{Iván Bautista Alvarado$^1$ y Andrés López González$^1$}
    \address{$^1$ Facultad de Ciencias, Universidad Nacional Autónoma de México, M\'exico CDMX.}
    \begin{abstract}
    Registro experimental sobre la redistribución de carga eléctrica en materiales. Experimentos: análisis de carga por fricción en electroscopio, y separación de bolas de unicel inducidas con combinaciones material-tela. 
    \end{abstract}
    \begin{keyword}
    Electromagnetismo; Electrostática; Fuerza de Coulomb; Redistribución de carga
    \end{keyword}
\end{frontmatter}

\section{Introducción}
Experimento de prácticas de electrostática para analizar la carga eléctrica de objetos comunes, comparable a la propiedad de masa en física gravitacional. Cargas aplicadas mediante combinación de barras y telas, registro de fuerza y deflexión en electroscopio.

\section{Metodología}
\subsection{Etapa 1: Medición de Carga por Fricción}
Materiales: barras (PVC, latón, madera, etc.), telas (licra, polar, terciopelo, etc.), electroscopio y péndulo de unicel. \\
Proceso: Frotado de cada barra con tela durante 30 segundos; acercamiento al electroscopio y péndulo de unicel, sin contacto. Se registran ángulos de deflexión.

\subsection{Etapa 2: Inducción de Voltaje en Electroscopio}
Conexión de fuente de alto voltaje (kV) al electroscopio. Ajuste gradual hasta registrar ángulos de deflexión de 5° a 55°. \\
Materiales: electroscopio, fuente de alto voltaje.

\subsection{Etapa 3: Separación de Bolas de Unicel por Inducción}
Materiales: montaje de péndulo doble con bolas de unicel, selección de tres combinaciones de material-tela con mayor carga registrada. \\
Proceso: Frotado de barras seleccionadas; acercamiento a bolas de unicel suspendidas, registrando el ángulo de separación inducido.

\section{Registro de Resultados}
\subsection{Resultados Cualitativos y Cuantitativos}
Escala cualitativa inicial (1-10.5) ajustada para reflejar mayor intensidad de deflexión observada. Tabla de ángulos registrados en electroscopio para cada combinación material-tela; combinaciones con mayor carga resaltadas.

\subsection{Inducción de Voltaje en Electroscopio}
Medición de ángulos de deflexión frente a voltaje aplicado (kV). Observación de relación no lineal, con mejor ajuste cuadrático entre ángulo y raíz de voltaje. \\
Tabla de datos: voltaje aplicado vs. ángulo de deflexión.

\subsection{Carga en Bolas de Unicel}
Cálculo de carga en bolas de unicel mediante deflexión inducida; resultados de las tres combinaciones seleccionadas con valores e incertidumbres presentados en gráfico.

\section{Observaciones y Notas}
Ajustes: escala cualitativa revisada tras observación de mayores deflexiones en combinación lana-vidrio. El análisis cualitativo resultó menos preciso comparado con la medición directa del electroscopio. \\
\textbf{Etapa 2}: Observación de relación no lineal entre voltaje aplicado y ángulo, con incertidumbre en cada registro.
\textbf{Etapa 3}: Los ángulos de separación entre bolas de unicel reflejan cargas inducidas proporcionalmente a la afinidad del material por carga de fricción.

\end{document}
